% =============================================================================
% Integrante: Enrique A Gonzalez Moreira
% Fragmento puro: se incluye desde preguntas/p4-inversion-dependencias/solucion.tex
% mediante % =============================================================================
% Integrante: Enrique A Gonzalez Moreira
% Fragmento puro: se incluye desde preguntas/p4-inversion-dependencias/solucion.tex
% mediante % =============================================================================
% Integrante: Enrique A Gonzalez Moreira
% Fragmento puro: se incluye desde preguntas/p4-inversion-dependencias/solucion.tex
% mediante % =============================================================================
% Integrante: Enrique A Gonzalez Moreira
% Fragmento puro: se incluye desde preguntas/p4-inversion-dependencias/solucion.tex
% mediante \input{investigacion/p4-dip-interfaces/notas} (ya presente en ese
% archivo: sin este fragmento, SemArq1_TeamBales.tex no compila).
% Fuente de trabajo: notas.md (misma carpeta).
% Citas: páginas impresas de Martin (2017); Cockburn (2005) y Palermo (2008) no
% tienen paginación y se citan por sección o parte. Todas las claves están en
% core.bib (martin2017, cockburn2005, palermo2008): fuentes.bib no declara nada.
% Comillas: se usan comillas tipográficas literales, no \enquote, para no
% depender de csquotes (no está en la lista de paquetes de preamble.tex, #13).
% =============================================================================

% Grupo local: \emergencystretch evita líneas desbordadas sin afectar a otros capítulos.
\begingroup\emergencystretch=3em \hbadness=2000
\section*{Notas de investigación: DIP, Inversión de Dependencias y diseño de interfaces en Clean Architecture}

\section{Síntesis}

El Principio de Inversión de Dependencias (DIP) sostiene que los sistemas más flexibles son
aquellos en los que las dependencias del código fuente se refieren solo a abstracciones, no a
concreciones \parencite[87]{martin2017}. La consecuencia práctica es que una política de alto
nivel —una regla de negocio— no debe nombrar un detalle de bajo nivel: ni una base de datos, ni
un servicio externo, ni una librería de transporte. Ambos, política y detalle, dependen de una
abstracción declarada en el núcleo. Martin añade el criterio que convierte el principio en
decisión de diseño: toda modificación de una interfaz abstracta obliga a modificar sus
implementaciones, pero no al contrario, de modo que las interfaces son menos volátiles que las
implementaciones \parencite[88]{martin2017}; de ahí las reglas operativas de no referirse a
clases concretas volátiles, no sobreescribir funciones concretas y recurrir a factorías
abstractas para la creación de objetos \parencite[89]{martin2017}.

La Regla de Dependencia formaliza el principio a escala arquitectónica: las dependencias del
código fuente deben apuntar solo hacia adentro, hacia políticas de mayor nivel, y nada de un
círculo interno puede conocer nada de un círculo externo \parencite[203]{martin2017}. Martin
presenta esa regla como denominador común de un conjunto de arquitecturas —Hexagonal de
Cockburn, DCI, BCE— que persiguen la misma separación de intereses y producen sistemas
independientes de frameworks, de la interfaz de usuario, de la base de datos y de cualquier
agencia externa, y con reglas de negocio probables sin ninguno de esos elementos
\parencite[202]{martin2017}. En la Hexagonal, el núcleo declara los \emph{puertos} y la
tecnología aporta los \emph{adaptadores} que los implementan \parencite{cockburn2005}; en la
Onion, el Modelo de Dominio ocupa el centro y las interfaces de repositorio viven en el núcleo
mientras el comportamiento que toca la base de datos queda fuera \parencite{palermo2008}.

El mecanismo concreto es siempre el mismo: el detalle implementa la abstracción que el núcleo
declara, y la traducción entre el contrato del núcleo y la tecnología externa ocurre en la capa
de adaptadores de interfaz \parencite[205]{martin2017}. Conviene no confundirlo con la inversión
del flujo de control: el núcleo sigue llamando a la validación o a la consulta; lo que se
invierte es la dirección de la dependencia del código fuente, que en el límite cruzado pasa a
apuntar desde el detalle hacia la abstracción \parencite[87, 91]{martin2017}. Ese es el punto de
apoyo de la testabilidad: con el puerto en el núcleo, la prueba sustituye el adaptador real por
un doble controlado y ejercita la regla de negocio sin red, sin servidor y sin base de datos
\parencite[202]{martin2017}.

\section{Conceptos clave}

\begin{table}[htbp]
  \centering
  \footnotesize
  \renewcommand{\arraystretch}{1.2}
  \begin{tabular}{@{}p{0.22\linewidth}p{0.46\linewidth}p{0.24\linewidth}@{}}
    \toprule
    \textbf{Concepto} & \textbf{Definición} & \textbf{Fuente} \\
    \midrule
    \raggedright Principio de Inversión de Dependencias (DIP)
      & \raggedright Los sistemas más flexibles son los que dependen solo de abstracciones, no
        de concreciones; la política de alto nivel no nombra el detalle
      & \raggedright \textcite[87]{martin2017} \tabularnewline
    \raggedright Volatilidad de interfaces frente a implementaciones
      & \raggedright Cambiar una interfaz obliga a cambiar sus implementaciones, pero no al
        contrario: la interfaz es el elemento más estable del que depender
      & \raggedright \textcite[88]{martin2017} \tabularnewline
    \raggedright Reglas operativas del DIP
      & \raggedright No referirse a clases concretas volátiles, no sobreescribir funciones
        concretas y usar factorías abstractas para la creación
      & \raggedright \textcite[89]{martin2017} \tabularnewline
    \raggedright Regla de Dependencia
      & \raggedright Las dependencias del código fuente apuntan solo hacia adentro; nada de un
        círculo interno conoce nada de un círculo externo
      & \raggedright \textcite[203]{martin2017} \tabularnewline
    \raggedright Características comunes de Clean, Hexagonal, DCI y BCE
      & \raggedright Independencia de frameworks, de la interfaz de usuario, de la base de datos
        y de agencias externas; reglas de negocio probables en aislamiento
      & \raggedright \textcite[202]{martin2017} \tabularnewline
    \raggedright Puertos y adaptadores
      & \raggedright El núcleo declara el puerto y la tecnología aporta el adaptador que lo
        implementa, de modo que el núcleo ignora el dispositivo concreto
      & \raggedright \textcite{cockburn2005} \tabularnewline
    \raggedright Capa de adaptadores de interfaz
      & \raggedright Convierte los datos entre el formato del caso de uso y el de la tecnología
        externa; es el lugar natural del detalle
      & \raggedright \textcite[205]{martin2017} \tabularnewline
    \raggedright Dominio al centro (Onion)
      & \raggedright El Modelo de Dominio no se acopla más que a sí mismo; las interfaces de
        repositorio van en el núcleo y el acceso a datos fuera
      & \raggedright \textcite{palermo2008} \tabularnewline
    \raggedright Flujo de control frente a dependencia de código fuente
      & \raggedright El núcleo sigue invocando la operación; lo que se invierte en el límite es
        la dirección de la dependencia del código fuente
      & \raggedright \textcite[87, 91]{martin2017} \tabularnewline
    \raggedright Testabilidad mediante dobles de prueba
      & \raggedright Con el puerto en el núcleo, la prueba sustituye el adaptador real y ejercita
        la regla de negocio sin interfaz, base de datos ni servidor
      & \raggedright \textcite[202]{martin2017} \tabularnewline
    \bottomrule
  \end{tabular}
  \caption{Conceptos clave de la investigación INV-P4.}
  \label{tab:inv-p4-conceptos}
\end{table}

\section{Puntos de conexión con las preguntas}

\begin{itemize}
  \item \textbf{P4, inciso a):} la dependencia directa \texttt{OrderProcessor}
        $\rightarrow$ \texttt{AddressValidatorService} hace que una política de alto nivel
        nombre un detalle volátil, lo que contraviene el DIP \parencite[87]{martin2017} y la
        Regla de Dependencia \parencite[203]{martin2017}, y rompe la relación puerto-adaptador
        \parencite{cockburn2005}.
  \item \textbf{P4, inciso b):} el rediseño consiste en declarar \texttt{IAddressValidator} en
        el núcleo y dejar el adaptador HTTP fuera, en la capa de adaptadores de interfaz
        \parencite[205]{martin2017}, por ser el elemento volátil \parencite[88]{martin2017}.
  \item \textbf{P4, inciso c):} la prueba del núcleo sin HTTP se sostiene en la testabilidad de
        estas arquitecturas \parencite[202]{martin2017}, y el cambio de proveedor queda
        confinado al adaptador porque la dependencia apunta a la interfaz, el elemento más
        estable \parencite[88]{martin2017}.
  \item \textbf{P1, P2 y P6:} la Regla de Dependencia y la simetría puerto-adaptador son el
        mismo criterio que P1 y P2 contrastan con la N-Capas, y el argumento estructural que P6
        opone a las optimizaciones de capas abiertas.
\end{itemize}


 (ya presente en ese
% archivo: sin este fragmento, SemArq1_TeamBales.tex no compila).
% Fuente de trabajo: notas.md (misma carpeta).
% Citas: páginas impresas de Martin (2017); Cockburn (2005) y Palermo (2008) no
% tienen paginación y se citan por sección o parte. Todas las claves están en
% core.bib (martin2017, cockburn2005, palermo2008): fuentes.bib no declara nada.
% Comillas: se usan comillas tipográficas literales, no \enquote, para no
% depender de csquotes (no está en la lista de paquetes de preamble.tex, #13).
% =============================================================================

% Grupo local: \emergencystretch evita líneas desbordadas sin afectar a otros capítulos.
\begingroup\emergencystretch=3em \hbadness=2000
\section*{Notas de investigación: DIP, Inversión de Dependencias y diseño de interfaces en Clean Architecture}

\section{Síntesis}

El Principio de Inversión de Dependencias (DIP) sostiene que los sistemas más flexibles son
aquellos en los que las dependencias del código fuente se refieren solo a abstracciones, no a
concreciones \parencite[87]{martin2017}. La consecuencia práctica es que una política de alto
nivel —una regla de negocio— no debe nombrar un detalle de bajo nivel: ni una base de datos, ni
un servicio externo, ni una librería de transporte. Ambos, política y detalle, dependen de una
abstracción declarada en el núcleo. Martin añade el criterio que convierte el principio en
decisión de diseño: toda modificación de una interfaz abstracta obliga a modificar sus
implementaciones, pero no al contrario, de modo que las interfaces son menos volátiles que las
implementaciones \parencite[88]{martin2017}; de ahí las reglas operativas de no referirse a
clases concretas volátiles, no sobreescribir funciones concretas y recurrir a factorías
abstractas para la creación de objetos \parencite[89]{martin2017}.

La Regla de Dependencia formaliza el principio a escala arquitectónica: las dependencias del
código fuente deben apuntar solo hacia adentro, hacia políticas de mayor nivel, y nada de un
círculo interno puede conocer nada de un círculo externo \parencite[203]{martin2017}. Martin
presenta esa regla como denominador común de un conjunto de arquitecturas —Hexagonal de
Cockburn, DCI, BCE— que persiguen la misma separación de intereses y producen sistemas
independientes de frameworks, de la interfaz de usuario, de la base de datos y de cualquier
agencia externa, y con reglas de negocio probables sin ninguno de esos elementos
\parencite[202]{martin2017}. En la Hexagonal, el núcleo declara los \emph{puertos} y la
tecnología aporta los \emph{adaptadores} que los implementan \parencite{cockburn2005}; en la
Onion, el Modelo de Dominio ocupa el centro y las interfaces de repositorio viven en el núcleo
mientras el comportamiento que toca la base de datos queda fuera \parencite{palermo2008}.

El mecanismo concreto es siempre el mismo: el detalle implementa la abstracción que el núcleo
declara, y la traducción entre el contrato del núcleo y la tecnología externa ocurre en la capa
de adaptadores de interfaz \parencite[205]{martin2017}. Conviene no confundirlo con la inversión
del flujo de control: el núcleo sigue llamando a la validación o a la consulta; lo que se
invierte es la dirección de la dependencia del código fuente, que en el límite cruzado pasa a
apuntar desde el detalle hacia la abstracción \parencite[87, 91]{martin2017}. Ese es el punto de
apoyo de la testabilidad: con el puerto en el núcleo, la prueba sustituye el adaptador real por
un doble controlado y ejercita la regla de negocio sin red, sin servidor y sin base de datos
\parencite[202]{martin2017}.

\section{Conceptos clave}

\begin{table}[htbp]
  \centering
  \footnotesize
  \renewcommand{\arraystretch}{1.2}
  \begin{tabular}{@{}p{0.22\linewidth}p{0.46\linewidth}p{0.24\linewidth}@{}}
    \toprule
    \textbf{Concepto} & \textbf{Definición} & \textbf{Fuente} \\
    \midrule
    \raggedright Principio de Inversión de Dependencias (DIP)
      & \raggedright Los sistemas más flexibles son los que dependen solo de abstracciones, no
        de concreciones; la política de alto nivel no nombra el detalle
      & \raggedright \textcite[87]{martin2017} \tabularnewline
    \raggedright Volatilidad de interfaces frente a implementaciones
      & \raggedright Cambiar una interfaz obliga a cambiar sus implementaciones, pero no al
        contrario: la interfaz es el elemento más estable del que depender
      & \raggedright \textcite[88]{martin2017} \tabularnewline
    \raggedright Reglas operativas del DIP
      & \raggedright No referirse a clases concretas volátiles, no sobreescribir funciones
        concretas y usar factorías abstractas para la creación
      & \raggedright \textcite[89]{martin2017} \tabularnewline
    \raggedright Regla de Dependencia
      & \raggedright Las dependencias del código fuente apuntan solo hacia adentro; nada de un
        círculo interno conoce nada de un círculo externo
      & \raggedright \textcite[203]{martin2017} \tabularnewline
    \raggedright Características comunes de Clean, Hexagonal, DCI y BCE
      & \raggedright Independencia de frameworks, de la interfaz de usuario, de la base de datos
        y de agencias externas; reglas de negocio probables en aislamiento
      & \raggedright \textcite[202]{martin2017} \tabularnewline
    \raggedright Puertos y adaptadores
      & \raggedright El núcleo declara el puerto y la tecnología aporta el adaptador que lo
        implementa, de modo que el núcleo ignora el dispositivo concreto
      & \raggedright \textcite{cockburn2005} \tabularnewline
    \raggedright Capa de adaptadores de interfaz
      & \raggedright Convierte los datos entre el formato del caso de uso y el de la tecnología
        externa; es el lugar natural del detalle
      & \raggedright \textcite[205]{martin2017} \tabularnewline
    \raggedright Dominio al centro (Onion)
      & \raggedright El Modelo de Dominio no se acopla más que a sí mismo; las interfaces de
        repositorio van en el núcleo y el acceso a datos fuera
      & \raggedright \textcite{palermo2008} \tabularnewline
    \raggedright Flujo de control frente a dependencia de código fuente
      & \raggedright El núcleo sigue invocando la operación; lo que se invierte en el límite es
        la dirección de la dependencia del código fuente
      & \raggedright \textcite[87, 91]{martin2017} \tabularnewline
    \raggedright Testabilidad mediante dobles de prueba
      & \raggedright Con el puerto en el núcleo, la prueba sustituye el adaptador real y ejercita
        la regla de negocio sin interfaz, base de datos ni servidor
      & \raggedright \textcite[202]{martin2017} \tabularnewline
    \bottomrule
  \end{tabular}
  \caption{Conceptos clave de la investigación INV-P4.}
  \label{tab:inv-p4-conceptos}
\end{table}

\section{Puntos de conexión con las preguntas}

\begin{itemize}
  \item \textbf{P4, inciso a):} la dependencia directa \texttt{OrderProcessor}
        $\rightarrow$ \texttt{AddressValidatorService} hace que una política de alto nivel
        nombre un detalle volátil, lo que contraviene el DIP \parencite[87]{martin2017} y la
        Regla de Dependencia \parencite[203]{martin2017}, y rompe la relación puerto-adaptador
        \parencite{cockburn2005}.
  \item \textbf{P4, inciso b):} el rediseño consiste en declarar \texttt{IAddressValidator} en
        el núcleo y dejar el adaptador HTTP fuera, en la capa de adaptadores de interfaz
        \parencite[205]{martin2017}, por ser el elemento volátil \parencite[88]{martin2017}.
  \item \textbf{P4, inciso c):} la prueba del núcleo sin HTTP se sostiene en la testabilidad de
        estas arquitecturas \parencite[202]{martin2017}, y el cambio de proveedor queda
        confinado al adaptador porque la dependencia apunta a la interfaz, el elemento más
        estable \parencite[88]{martin2017}.
  \item \textbf{P1, P2 y P6:} la Regla de Dependencia y la simetría puerto-adaptador son el
        mismo criterio que P1 y P2 contrastan con la N-Capas, y el argumento estructural que P6
        opone a las optimizaciones de capas abiertas.
\end{itemize}


 (ya presente en ese
% archivo: sin este fragmento, SemArq1_TeamBales.tex no compila).
% Fuente de trabajo: notas.md (misma carpeta).
% Citas: páginas impresas de Martin (2017); Cockburn (2005) y Palermo (2008) no
% tienen paginación y se citan por sección o parte. Todas las claves están en
% core.bib (martin2017, cockburn2005, palermo2008): fuentes.bib no declara nada.
% Comillas: se usan comillas tipográficas literales, no \enquote, para no
% depender de csquotes (no está en la lista de paquetes de preamble.tex, #13).
% =============================================================================

% Grupo local: \emergencystretch evita líneas desbordadas sin afectar a otros capítulos.
\begingroup\emergencystretch=3em \hbadness=2000
\section*{Notas de investigación: DIP, Inversión de Dependencias y diseño de interfaces en Clean Architecture}

\section{Síntesis}

El Principio de Inversión de Dependencias (DIP) sostiene que los sistemas más flexibles son
aquellos en los que las dependencias del código fuente se refieren solo a abstracciones, no a
concreciones \parencite[87]{martin2017}. La consecuencia práctica es que una política de alto
nivel —una regla de negocio— no debe nombrar un detalle de bajo nivel: ni una base de datos, ni
un servicio externo, ni una librería de transporte. Ambos, política y detalle, dependen de una
abstracción declarada en el núcleo. Martin añade el criterio que convierte el principio en
decisión de diseño: toda modificación de una interfaz abstracta obliga a modificar sus
implementaciones, pero no al contrario, de modo que las interfaces son menos volátiles que las
implementaciones \parencite[88]{martin2017}; de ahí las reglas operativas de no referirse a
clases concretas volátiles, no sobreescribir funciones concretas y recurrir a factorías
abstractas para la creación de objetos \parencite[89]{martin2017}.

La Regla de Dependencia formaliza el principio a escala arquitectónica: las dependencias del
código fuente deben apuntar solo hacia adentro, hacia políticas de mayor nivel, y nada de un
círculo interno puede conocer nada de un círculo externo \parencite[203]{martin2017}. Martin
presenta esa regla como denominador común de un conjunto de arquitecturas —Hexagonal de
Cockburn, DCI, BCE— que persiguen la misma separación de intereses y producen sistemas
independientes de frameworks, de la interfaz de usuario, de la base de datos y de cualquier
agencia externa, y con reglas de negocio probables sin ninguno de esos elementos
\parencite[202]{martin2017}. En la Hexagonal, el núcleo declara los \emph{puertos} y la
tecnología aporta los \emph{adaptadores} que los implementan \parencite{cockburn2005}; en la
Onion, el Modelo de Dominio ocupa el centro y las interfaces de repositorio viven en el núcleo
mientras el comportamiento que toca la base de datos queda fuera \parencite{palermo2008}.

El mecanismo concreto es siempre el mismo: el detalle implementa la abstracción que el núcleo
declara, y la traducción entre el contrato del núcleo y la tecnología externa ocurre en la capa
de adaptadores de interfaz \parencite[205]{martin2017}. Conviene no confundirlo con la inversión
del flujo de control: el núcleo sigue llamando a la validación o a la consulta; lo que se
invierte es la dirección de la dependencia del código fuente, que en el límite cruzado pasa a
apuntar desde el detalle hacia la abstracción \parencite[87, 91]{martin2017}. Ese es el punto de
apoyo de la testabilidad: con el puerto en el núcleo, la prueba sustituye el adaptador real por
un doble controlado y ejercita la regla de negocio sin red, sin servidor y sin base de datos
\parencite[202]{martin2017}.

\section{Conceptos clave}

\begin{table}[htbp]
  \centering
  \footnotesize
  \renewcommand{\arraystretch}{1.2}
  \begin{tabular}{@{}p{0.22\linewidth}p{0.46\linewidth}p{0.24\linewidth}@{}}
    \toprule
    \textbf{Concepto} & \textbf{Definición} & \textbf{Fuente} \\
    \midrule
    \raggedright Principio de Inversión de Dependencias (DIP)
      & \raggedright Los sistemas más flexibles son los que dependen solo de abstracciones, no
        de concreciones; la política de alto nivel no nombra el detalle
      & \raggedright \textcite[87]{martin2017} \tabularnewline
    \raggedright Volatilidad de interfaces frente a implementaciones
      & \raggedright Cambiar una interfaz obliga a cambiar sus implementaciones, pero no al
        contrario: la interfaz es el elemento más estable del que depender
      & \raggedright \textcite[88]{martin2017} \tabularnewline
    \raggedright Reglas operativas del DIP
      & \raggedright No referirse a clases concretas volátiles, no sobreescribir funciones
        concretas y usar factorías abstractas para la creación
      & \raggedright \textcite[89]{martin2017} \tabularnewline
    \raggedright Regla de Dependencia
      & \raggedright Las dependencias del código fuente apuntan solo hacia adentro; nada de un
        círculo interno conoce nada de un círculo externo
      & \raggedright \textcite[203]{martin2017} \tabularnewline
    \raggedright Características comunes de Clean, Hexagonal, DCI y BCE
      & \raggedright Independencia de frameworks, de la interfaz de usuario, de la base de datos
        y de agencias externas; reglas de negocio probables en aislamiento
      & \raggedright \textcite[202]{martin2017} \tabularnewline
    \raggedright Puertos y adaptadores
      & \raggedright El núcleo declara el puerto y la tecnología aporta el adaptador que lo
        implementa, de modo que el núcleo ignora el dispositivo concreto
      & \raggedright \textcite{cockburn2005} \tabularnewline
    \raggedright Capa de adaptadores de interfaz
      & \raggedright Convierte los datos entre el formato del caso de uso y el de la tecnología
        externa; es el lugar natural del detalle
      & \raggedright \textcite[205]{martin2017} \tabularnewline
    \raggedright Dominio al centro (Onion)
      & \raggedright El Modelo de Dominio no se acopla más que a sí mismo; las interfaces de
        repositorio van en el núcleo y el acceso a datos fuera
      & \raggedright \textcite{palermo2008} \tabularnewline
    \raggedright Flujo de control frente a dependencia de código fuente
      & \raggedright El núcleo sigue invocando la operación; lo que se invierte en el límite es
        la dirección de la dependencia del código fuente
      & \raggedright \textcite[87, 91]{martin2017} \tabularnewline
    \raggedright Testabilidad mediante dobles de prueba
      & \raggedright Con el puerto en el núcleo, la prueba sustituye el adaptador real y ejercita
        la regla de negocio sin interfaz, base de datos ni servidor
      & \raggedright \textcite[202]{martin2017} \tabularnewline
    \bottomrule
  \end{tabular}
  \caption{Conceptos clave de la investigación INV-P4.}
  \label{tab:inv-p4-conceptos}
\end{table}

\section{Puntos de conexión con las preguntas}

\begin{itemize}
  \item \textbf{P4, inciso a):} la dependencia directa \texttt{OrderProcessor}
        $\rightarrow$ \texttt{AddressValidatorService} hace que una política de alto nivel
        nombre un detalle volátil, lo que contraviene el DIP \parencite[87]{martin2017} y la
        Regla de Dependencia \parencite[203]{martin2017}, y rompe la relación puerto-adaptador
        \parencite{cockburn2005}.
  \item \textbf{P4, inciso b):} el rediseño consiste en declarar \texttt{IAddressValidator} en
        el núcleo y dejar el adaptador HTTP fuera, en la capa de adaptadores de interfaz
        \parencite[205]{martin2017}, por ser el elemento volátil \parencite[88]{martin2017}.
  \item \textbf{P4, inciso c):} la prueba del núcleo sin HTTP se sostiene en la testabilidad de
        estas arquitecturas \parencite[202]{martin2017}, y el cambio de proveedor queda
        confinado al adaptador porque la dependencia apunta a la interfaz, el elemento más
        estable \parencite[88]{martin2017}.
  \item \textbf{P1, P2 y P6:} la Regla de Dependencia y la simetría puerto-adaptador son el
        mismo criterio que P1 y P2 contrastan con la N-Capas, y el argumento estructural que P6
        opone a las optimizaciones de capas abiertas.
\end{itemize}


 (ya presente en ese
% archivo: sin este fragmento, SemArq1_TeamBales.tex no compila).
% Fuente de trabajo: notas.md (misma carpeta).
% Citas: páginas impresas de Martin (2017); Cockburn (2005) y Palermo (2008) no
% tienen paginación y se citan por sección o parte. Todas las claves están en
% core.bib (martin2017, cockburn2005, palermo2008): fuentes.bib no declara nada.
% Comillas: se usan comillas tipográficas literales, no \enquote, para no
% depender de csquotes (no está en la lista de paquetes de preamble.tex, #13).
% =============================================================================

% Grupo local: \emergencystretch evita líneas desbordadas sin afectar a otros capítulos.
\begingroup\emergencystretch=3em \hbadness=2000
\section*{Notas de investigación: DIP, Inversión de Dependencias y diseño de interfaces en Clean Architecture}

\section{Síntesis}

El Principio de Inversión de Dependencias (DIP) sostiene que los sistemas más flexibles son
aquellos en los que las dependencias del código fuente se refieren solo a abstracciones, no a
concreciones \parencite[87]{martin2017}. La consecuencia práctica es que una política de alto
nivel —una regla de negocio— no debe nombrar un detalle de bajo nivel: ni una base de datos, ni
un servicio externo, ni una librería de transporte. Ambos, política y detalle, dependen de una
abstracción declarada en el núcleo. Martin añade el criterio que convierte el principio en
decisión de diseño: toda modificación de una interfaz abstracta obliga a modificar sus
implementaciones, pero no al contrario, de modo que las interfaces son menos volátiles que las
implementaciones \parencite[88]{martin2017}; de ahí las reglas operativas de no referirse a
clases concretas volátiles, no sobreescribir funciones concretas y recurrir a factorías
abstractas para la creación de objetos \parencite[89]{martin2017}.

La Regla de Dependencia formaliza el principio a escala arquitectónica: las dependencias del
código fuente deben apuntar solo hacia adentro, hacia políticas de mayor nivel, y nada de un
círculo interno puede conocer nada de un círculo externo \parencite[203]{martin2017}. Martin
presenta esa regla como denominador común de un conjunto de arquitecturas —Hexagonal de
Cockburn, DCI, BCE— que persiguen la misma separación de intereses y producen sistemas
independientes de frameworks, de la interfaz de usuario, de la base de datos y de cualquier
agencia externa, y con reglas de negocio probables sin ninguno de esos elementos
\parencite[202]{martin2017}. En la Hexagonal, el núcleo declara los \emph{puertos} y la
tecnología aporta los \emph{adaptadores} que los implementan \parencite{cockburn2005}; en la
Onion, el Modelo de Dominio ocupa el centro y las interfaces de repositorio viven en el núcleo
mientras el comportamiento que toca la base de datos queda fuera \parencite{palermo2008}.

El mecanismo concreto es siempre el mismo: el detalle implementa la abstracción que el núcleo
declara, y la traducción entre el contrato del núcleo y la tecnología externa ocurre en la capa
de adaptadores de interfaz \parencite[205]{martin2017}. Conviene no confundirlo con la inversión
del flujo de control: el núcleo sigue llamando a la validación o a la consulta; lo que se
invierte es la dirección de la dependencia del código fuente, que en el límite cruzado pasa a
apuntar desde el detalle hacia la abstracción \parencite[87, 91]{martin2017}. Ese es el punto de
apoyo de la testabilidad: con el puerto en el núcleo, la prueba sustituye el adaptador real por
un doble controlado y ejercita la regla de negocio sin red, sin servidor y sin base de datos
\parencite[202]{martin2017}.

\section{Conceptos clave}

\begin{table}[htbp]
  \centering
  \footnotesize
  \renewcommand{\arraystretch}{1.2}
  \begin{tabular}{@{}p{0.22\linewidth}p{0.46\linewidth}p{0.24\linewidth}@{}}
    \toprule
    \textbf{Concepto} & \textbf{Definición} & \textbf{Fuente} \\
    \midrule
    \raggedright Principio de Inversión de Dependencias (DIP)
      & \raggedright Los sistemas más flexibles son los que dependen solo de abstracciones, no
        de concreciones; la política de alto nivel no nombra el detalle
      & \raggedright \textcite[87]{martin2017} \tabularnewline
    \raggedright Volatilidad de interfaces frente a implementaciones
      & \raggedright Cambiar una interfaz obliga a cambiar sus implementaciones, pero no al
        contrario: la interfaz es el elemento más estable del que depender
      & \raggedright \textcite[88]{martin2017} \tabularnewline
    \raggedright Reglas operativas del DIP
      & \raggedright No referirse a clases concretas volátiles, no sobreescribir funciones
        concretas y usar factorías abstractas para la creación
      & \raggedright \textcite[89]{martin2017} \tabularnewline
    \raggedright Regla de Dependencia
      & \raggedright Las dependencias del código fuente apuntan solo hacia adentro; nada de un
        círculo interno conoce nada de un círculo externo
      & \raggedright \textcite[203]{martin2017} \tabularnewline
    \raggedright Características comunes de Clean, Hexagonal, DCI y BCE
      & \raggedright Independencia de frameworks, de la interfaz de usuario, de la base de datos
        y de agencias externas; reglas de negocio probables en aislamiento
      & \raggedright \textcite[202]{martin2017} \tabularnewline
    \raggedright Puertos y adaptadores
      & \raggedright El núcleo declara el puerto y la tecnología aporta el adaptador que lo
        implementa, de modo que el núcleo ignora el dispositivo concreto
      & \raggedright \textcite{cockburn2005} \tabularnewline
    \raggedright Capa de adaptadores de interfaz
      & \raggedright Convierte los datos entre el formato del caso de uso y el de la tecnología
        externa; es el lugar natural del detalle
      & \raggedright \textcite[205]{martin2017} \tabularnewline
    \raggedright Dominio al centro (Onion)
      & \raggedright El Modelo de Dominio no se acopla más que a sí mismo; las interfaces de
        repositorio van en el núcleo y el acceso a datos fuera
      & \raggedright \textcite{palermo2008} \tabularnewline
    \raggedright Flujo de control frente a dependencia de código fuente
      & \raggedright El núcleo sigue invocando la operación; lo que se invierte en el límite es
        la dirección de la dependencia del código fuente
      & \raggedright \textcite[87, 91]{martin2017} \tabularnewline
    \raggedright Testabilidad mediante dobles de prueba
      & \raggedright Con el puerto en el núcleo, la prueba sustituye el adaptador real y ejercita
        la regla de negocio sin interfaz, base de datos ni servidor
      & \raggedright \textcite[202]{martin2017} \tabularnewline
    \bottomrule
  \end{tabular}
  \caption{Conceptos clave de la investigación INV-P4.}
  \label{tab:inv-p4-conceptos}
\end{table}

\section{Puntos de conexión con las preguntas}

\begin{itemize}
  \item \textbf{P4, inciso a):} la dependencia directa \texttt{OrderProcessor}
        $\rightarrow$ \texttt{AddressValidatorService} hace que una política de alto nivel
        nombre un detalle volátil, lo que contraviene el DIP \parencite[87]{martin2017} y la
        Regla de Dependencia \parencite[203]{martin2017}, y rompe la relación puerto-adaptador
        \parencite{cockburn2005}.
  \item \textbf{P4, inciso b):} el rediseño consiste en declarar \texttt{IAddressValidator} en
        el núcleo y dejar el adaptador HTTP fuera, en la capa de adaptadores de interfaz
        \parencite[205]{martin2017}, por ser el elemento volátil \parencite[88]{martin2017}.
  \item \textbf{P4, inciso c):} la prueba del núcleo sin HTTP se sostiene en la testabilidad de
        estas arquitecturas \parencite[202]{martin2017}, y el cambio de proveedor queda
        confinado al adaptador porque la dependencia apunta a la interfaz, el elemento más
        estable \parencite[88]{martin2017}.
  \item \textbf{P1, P2 y P6:} la Regla de Dependencia y la simetría puerto-adaptador son el
        mismo criterio que P1 y P2 contrastan con la N-Capas, y el argumento estructural que P6
        opone a las optimizaciones de capas abiertas.
\end{itemize}


